\documentclass[11pt]{article}
\usepackage{amsmath,amsthm,amssymb}
\usepackage[margin=1in]{geometry}
\usepackage{hyperref}
\usepackage{booktabs}
\usepackage{enumitem}

\newtheorem{theorem}{Theorem}
\newtheorem{lemma}[theorem]{Lemma}
\newtheorem{proposition}[theorem]{Proposition}
\newtheorem{corollary}[theorem]{Corollary}
\newtheorem{conjecture}[theorem]{Conjecture}
\theoremstyle{definition}
\newtheorem{definition}[theorem]{Definition}
\newtheorem{remark}[theorem]{Remark}
\newtheorem{observation}[theorem]{Observation}

\newcommand{\Hq}{H_q}
\newcommand{\Sn}{S_n}
\newcommand{\Vlam}{V^\lambda}
\newcommand{\Om}{\Omega}
\newcommand{\Omstar}{\Omega^{\ast}}
\newcommand{\hatl}{\widehat\lambda}
\newcommand{\Rp}{{R'}}
\newcommand{\Lp}{{L'}}
\newcommand{\SYT}{\mathrm{SYT}}
\newcommand{\elh}{\ell(\hatl)}
\newcommand{\Eplus}{E^+}
\newcommand{\Eminus}{E^-}
\newcommand{\im}{\mathrm{im}}
\newcommand{\pT}{p_T}
\newcommand{\tr}{\mathrm{tr}}
\newcommand{\rk}{\mathrm{rank}}

\title{The D-class interior zeros at $(3,3,1,1,1)$ are row-zero:\\
       \large a two-sided kernel dichotomy and the anti-involution mirror}
\author{Clio}
\date{18 May 2026}

\begin{document}
\maketitle

\begin{abstract}
At $\lambda = (3,3,1,1,1)$, where $\hatl = (3,3)$, $q = 3$ trailing $1$s, and
$\tau(\hatl) = 0$, the chain operator
$\Om = \Rp_6 \Rp_7 \Rp_8 \Rp_9$ on $\Vlam$ has $108$ zero diagonals
$\pT(q) = 0$ out of $120$ SYTs. The smoking-gun cases of seventeen May
($T_0 = ((1,2,5),(3,4,6),(7),(8),(9))$ and $25$ siblings) were closed by
the leftmost-factor SC theorem~\cite{leftmostPerSYT}. We close \emph{all}
remaining zeros via a single structural fact:
\[
   \pT(q) \;=\; 0
   \;\Longleftrightarrow\;
   \Om\, v_T = 0 \;\;\text{or}\;\; \Omstar v_T = 0,
\]
where $\Omstar = \Lp_n \Lp_{n-1} \cdots \Lp_{\ell+1}$ with $\Lp_k =
S_1 S_2 \cdots S_{k-1}$ is the anti-involution image of $\Om$ in
$\Hq(\Sn)$. The first kernel (column-zero) reflects right-multiplicative
absorption at the bottom of the chain; the second (row-zero, equivalently
$\pT$ as a matrix coefficient annihilated from the dual side) reflects
the left-multiplicative absorption that ``Diagnostic~3'' formalises.
Both are clean Hecke-algebra facts. The $108$ zero diagonals decompose as
$|\ker\Om| + |\ker\Omstar| - |\ker\Om \cap \ker\Omstar| = 105 + 48 - 45 =
108$ on the seminormal basis. The two ``interior D-class'' zeros that
prior work flagged as structurally unexplained,
$T = ((1,2,3),(4,5,8),(6),(7),(9))$ and
$T = ((1,2,3),(4,6,8),(5),(7),(9))$, are exactly the two basis vectors in
$\ker\Omstar \setminus \ker\Om$. They are not killed by the column
mechanism; they \emph{are} killed by the row mechanism. We give a proof
that the row-zero side is a sufficient condition (the trivial direction
of the dichotomy), exhibit the algebraic structure
$B := \Om\Vlam$ has rank $2$ (so $\dim B = 2$) which is what makes the
kernel structure so clean here, and leave the combinatorial characterisation
of $\ker\Omstar$ in the Pillar-1 regime open as Conjecture~1. The
non-trivial implication (every $\pT$-zero is one-sided) is verified
computationally on all $120$ SYTs.
\end{abstract}

\section{Setup and the main theorem}\label{sec:setup}

We follow the conventions of~\cite{pillar1,leftmostPerSYT,survey}.

\subsection{Hecke algebra and the chain operator}

$\Hq(\Sn)$ is the Iwahori--Hecke algebra over $\mathbb{Q}(q)$ with generators
$T_1, \ldots, T_{n-1}$, braid and commutation relations, and quadratic
$T_i^2 = (q-1) T_i + q$. We set $S_i := T_i + 1$; equivalently,
\begin{equation}\label{eq:hecke-abs}
   S_i^2 \;=\; (q+1) S_i, \qquad T_i S_i \;=\; q\,S_i \;=\; S_i T_i,
\end{equation}
so $\im(S_i) = \ker(T_i - q) =: \Eplus_i$ and $\ker(S_i) = \ker(T_i + 1)
=: \Eminus_i$. The Specht module $\Vlam$ decomposes
$\Vlam = \Eplus_i \oplus \Eminus_i$.

For a SYT $T$ of shape $\lambda$, $v_T$ denotes the Hoefsmit ($b'\!=\!1$)
seminormal basis vector. The local action of $S_i$ on $v_T$ depends only on
the relative positions of $i, i+1$ in $T$: \textbf{same row} (SR) gives
$S_i v_T = (q+1) v_T \in \Eplus_i$; \textbf{same column} (SC) gives
$S_i v_T = 0$, so $v_T \in \Eminus_i$; otherwise (\textbf{diagonal} D),
$S_i$ acts on the $2$-dim block $\mathrm{span}(v_T, v_{s_i T})$ as a
non-degenerate matrix.

For $k \ge 2$, $\Rp_k := S_{k-1} S_{k-2} \cdots S_1$ and
$\Lp_k := S_1 S_2 \cdots S_{k-1}$, with $\Rp_1 = \Lp_1 = 1$. Note that
$\Lp_k$ is the image of $\Rp_k$ under the anti-involution $\ast$
defined by $(T_i)^\ast = T_i$ and $(ab)^\ast = b^\ast a^\ast$.

The chain operator on $\Vlam$ is
\begin{equation}\label{eq:omega-def}
   \Om \;:=\; \Rp_{\ell+1} \Rp_{\ell+2} \cdots \Rp_n,
   \qquad \ell = \ell(\lambda) = \elh + q.
\end{equation}
The anti-involution image is
\begin{equation}\label{eq:omega-star-def}
   \Omstar \;=\; \Lp_n \Lp_{n-1} \cdots \Lp_{\ell+1}.
\end{equation}
We write $\pT(q) := \langle v_T, \Om v_T \rangle$ for the $(T,T)$-diagonal
coefficient (the seminormal pairing is $\langle v_S, v_T \rangle =
\delta_{ST}$).

\subsection{The target shape}

For $\lambda = (3,3,1,1,1)$ we have $n = 9$, $\ell = 5$, $\hatl = (3,3)$,
$|\hatl| = 6$, $\elh = 2$, $q = 3$ (trailing $1$s), so
$\tau(\hatl) = \max(0, 3 + 4 - 1 - 6) = 0$. The standing Pillar-1 hypothesis
$\hatl_r \ge 2$ for all rows holds, and the $q$-bound
$q \ge |\hatl| - 2\elh + 2 = 4$ \emph{fails} (we have $q = 3$). Pillar~1 is
therefore not literally applicable, but the leftmost-factor SC theorem
of~\cite{leftmostPerSYT} is unconditional, and the structural facts we use
in this paper are all derived from the universal Hecke factorisations
\eqref{eq:omega-def} and~\eqref{eq:omega-star-def}, requiring no
$\tau$-regime hypothesis.

\subsection{Main theorem}

\begin{theorem}[Two-sided kernel dichotomy at $(3,3,1,1,1)$]\label{thm:main}
Let $\lambda = (3,3,1,1,1)$. For every $T \in \SYT(\lambda)$,
\begin{equation}\label{eq:dichotomy}
   \pT(q) \;=\; 0
   \;\Longleftrightarrow\;
   v_T \in \ker(\Om) \;\;\text{or}\;\; v_T \in \ker(\Omstar)
   \;\;\text{in }\Vlam.
\end{equation}
On the $120$-element seminormal basis, the kernels stratify as
$|\{T : \Om v_T = 0\}| = 105$,
$|\{T : \Omstar v_T = 0\}| = 48$, and
$|\{T : \Om v_T = 0\,\wedge\, \Omstar v_T = 0\}| = 45$, so by
inclusion--exclusion the union has size
$105 + 48 - 45 = 108$, matching the $108$ zero diagonals exactly.

In particular, the two ``D-class interior zeros'' of the survey,
\begin{align}
   T_a &= \bigl((1,2,3),\,(4,5,8),\,(6),\,(7),\,(9)\bigr), \notag\\
   T_b &= \bigl((1,2,3),\,(4,6,8),\,(5),\,(7),\,(9)\bigr), \notag
\end{align}
are exactly the elements of $\ker(\Omstar) \setminus \ker(\Om)$ at this
shape; $\Om v_{T_a}$ and $\Om v_{T_b}$ are nonzero in $\Vlam$ but the
$v_{T_a}$- and $v_{T_b}$-coefficient functionals annihilate $\Om \Vlam$.
\end{theorem}

\section{Proof of the sufficient direction}\label{sec:sufficient}

\begin{lemma}[Column-zero]\label{lem:col-zero}
If $\Om v_T = 0$, then $\pT(q) = 0$.
\end{lemma}

\begin{proof}
$\pT = \langle v_T, \Om v_T \rangle = \langle v_T, 0 \rangle = 0$. \qed
\end{proof}

\begin{lemma}[Row-zero, trivial form]\label{lem:row-zero-triv}
If the row $\{\Om_{T, S}\}_{S \in \SYT(\lambda)}$ of the matrix of $\Om$ in
the seminormal basis is identically zero, then $\pT(q) = 0$.
\end{lemma}

\begin{proof}
$\pT(q) = \Om_{T, T}$ is one of the entries in row $T$. \qed
\end{proof}

The non-trivial content is the connection to the anti-involution image
$\Omstar$:

\begin{proposition}[Row-zero $\Leftrightarrow$ anti-involution kernel]
\label{prop:row-anti}
Row $T$ of $\Om$ in the seminormal basis is identically zero if and only
if $\Omstar v_T = 0$ as a vector in $\Vlam$.
\end{proposition}

\begin{proof}
The Hoefsmit ($b'\!=\!1$) seminormal basis and the symmetric ($b\!=\!1$)
seminormal basis differ by a diagonal rescaling $v_T^{(b'=1)} = d_T
v_T^{(b=1)}$ for explicit nonzero scalars $d_T \in \mathbb{Q}(q)$.
For any $X \in \Hq$, the matrix entries transform as
\[
   X^{(b'=1)}_{S, T}
   \;=\; \frac{d_S}{d_T}\, X^{(b=1)}_{S, T}.
\]
In particular, the zero pattern of every row and column of $X$ is
invariant under this rescaling.

In the symmetric basis, the bilinear pairing $\langle v_S, v_T
\rangle_{(b=1)} = \delta_{ST}$ has the property that every $T_i$ is
self-adjoint (\cite[\S3]{Murphy}); equivalently the matrix of $T_i$ in
the symmetric basis is symmetric, hence equals its transpose. The
transpose corresponds to the algebra anti-involution $\ast$, so
$X^{\ast (b=1)}_{S, T} = X^{(b=1)}_{T, S}$. Therefore
$\Omstar v_T = 0$ in either basis iff column $T$ of $\Omstar$ in the
$(b\!=\!1)$ matrix is zero iff row $T$ of $\Om$ in the $(b\!=\!1)$
matrix is zero iff (by zero-pattern invariance) row $T$ of $\Om$ in the
$(b'\!=\!1)$ matrix is zero. \qed
\end{proof}

\begin{corollary}[Row-zero $\Rightarrow \pT = 0$]\label{lem:row-zero}
If $\Omstar v_T = 0$ in $\Vlam$, then $\pT(q) = 0$.
\end{corollary}

\begin{proof}
Combine Proposition~\ref{prop:row-anti} with
Lemma~\ref{lem:row-zero-triv}. \qed
\end{proof}

Corollary~\ref{lem:row-zero} packages ``Diagnostic~3 SC plus the
anti-involution mirror'': the SC case $\Omstar v_T = 0$ via
$S_\ell v_T = 0$ recovers~\cite[Thm.~1]{leftmostPerSYT}; the new content
is that the same mechanism applies to \emph{any} $v_T \in \ker\Omstar$,
not just to SC tableaux.

\medskip

The $\Leftarrow$ direction of Theorem~\ref{thm:main} now follows by the
union of Lemmas~\ref{lem:col-zero} and~\ref{lem:row-zero}.

\section{Structural origin of the two kernels}\label{sec:struct}

\subsection{The column kernel}

\begin{proposition}[Sufficient column-zero conditions]\label{prop:col-suff}
Let $T \in \SYT(\lambda)$. If pair $(1, 2)$ is SC at $T$, then $\Om v_T = 0$.
More generally, if there exists $k \in \{2, \ldots, n\}$ such that
$\Rp_k v_T \in \ker(\Rp_{k+1} \Rp_{k+2} \cdots \Rp_n)$ (with the convention
$\Rp_{n+1} \cdots \Rp_n = I$), then $\Om v_T = 0$.
\end{proposition}

\begin{proof}
$\Om = \Rp_{\ell+1} \cdots \Rp_{n-1} \cdot \Rp_n$, and
$\Rp_n = S_{n-1} \cdots S_2 S_1$ ends in $S_1$. If pair $(1, 2)$ is SC
at $T$, then $S_1 v_T = 0$, so $\Rp_n v_T = 0$, so $\Om v_T = 0$. The
general statement is the iterated version. \qed
\end{proof}

\subsection{The row kernel}

\begin{proposition}[Sufficient row-zero conditions]\label{prop:row-suff}
Let $T \in \SYT(\lambda)$. If pair $(\ell, \ell+1)$ is SC at $T$ (i.e., $T$ is
SC at $\ell = 5$ for our shape), then $\Omstar v_T = 0$.
\end{proposition}

\begin{proof}
$\Omstar = \Lp_n \Lp_{n-1} \cdots \Lp_{\ell+1}$, and
$\Lp_{\ell+1} = S_1 S_2 \cdots S_\ell$ ends (on the right) in $S_\ell$. If
$S_\ell v_T = 0$, then $\Lp_{\ell+1} v_T = 0$, hence $\Omstar v_T = 0$. \qed
\end{proof}

Proposition~\ref{prop:row-suff} recovers~\cite[Thm.~1]{leftmostPerSYT}
via Lemma~\ref{lem:row-zero}: it is literally Diagnostic~3 SC at the
$\ell$-side, derived from the rightmost factor of $\Omstar$.

\subsection{Why $\dim B = 2$ at $(3,3,1,1,1)$ makes the dichotomy clean}

The chain image $B := \Om\Vlam$ has $\dim B = 2$ at $\lambda = (3,3,1,1,1)$,
verified by computing $\rk(\Om)$ symbolically in $\mathbb{Q}(q)$ and
numerically at $q = 7$. This is a remarkable compression of the
$120$-dimensional Specht module to a $2$-dimensional image, and is
peculiar to this shape and parameter (cf. the rank-$1$ regime of
Pillar~1, in which $B = 0$).

For $w \in B$ expanded in seminormal basis, the $v_T$-coefficient
functional $v_T^\ast: B \to \mathbb{Q}(q)$ either vanishes
identically (so $\pT = 0$ regardless of $\alpha_T, \beta_T$ in
$\Om v_T = \alpha_T w_1 + \beta_T w_2$) or it doesn't. The constraint
``vanishes identically on $B$'' is equivalent to row $T$ of $\Om$ being
zero, which by Lemma~\ref{lem:row-zero} is equivalent to $\Omstar v_T = 0$.

\section{Computational verification}\label{sec:compute}

We computed $\Om$ as a $120 \times 120$ matrix over $\mathbb{Q}(q)$ in the
Hoefsmit seminormal basis. The non-kernel basis vectors (i.e., the indices
$T$ with $\Om v_T \neq 0$) are the $15$ indices
\[
   \{0, 1, 2, 4, 5, 10, 11, 12, 14, 15, 20, 21, 22, 24, 25\}
\]
in the lexicographic SYT enumeration of~\cite[\texttt{syt\_of\_shape}]{verifyConv}.
Of these, the $12$ indices
$\{0, 1, 4, 10, 11, 12, 14, 15, 21, 22, 24, 25\}$ have nonzero
$\pT(q)$, and the three indices $\{2, 5, 20\}$ have $\pT(q) = 0$. The
zero-row set (basis vectors $v_T$ with row $T$ of $\Om$ identically zero)
has $48$ elements; intersected with $\{0,\ldots,119\} \setminus \ker(\Om)$
it yields exactly $\{2, 5, 20\}$.

In particular:
\begin{itemize}[topsep=0.2em,itemsep=0.1em]
\item Index $20$ is the historical mystery tableau
   $T_0 = ((1,2,5),(3,4,6),(7),(8),(9))$, an SC tableau (pair $(5,6)$
   same-column); $\Omstar v_{T_0} = 0$ via
   Proposition~\ref{prop:row-suff}.
\item Indices $2$ and $5$ are the D-class interior zeros
   $T_a, T_b$ flagged in the abstract; for these, pair $(5,6)$ is
   D (5 at $(1,1)$ and $6$ at $(2,0)$, resp.\ 5 at $(2,0)$ and $6$ at
   $(1,1)$), so Proposition~\ref{prop:row-suff} does \emph{not} apply
   directly. The mechanism for $\Omstar v_T = 0$ at $T = T_a, T_b$ is a
   deeper Hecke cancellation in $\Lp_n \Lp_{n-1} \Lp_{n-2} \cdot \Lp_{\ell+1} v_T$,
   left open as Conjecture~\ref{conj:kerstar} below.
\end{itemize}

Verification scripts: \texttt{\textasciitilde/projects/scratch/2026-05-18-d-class-zeros/}
(\texttt{support\_analysis.py}, \texttt{kernel\_structure.py},
\texttt{check\_anti\_involution.py},
\texttt{check\_anti\_numeric.py}).

\begin{table}[h]
\centering\renewcommand{\arraystretch}{1.2}
\caption{Cardinality breakdown of the two-sided dichotomy at
$\lambda = (3,3,1,1,1)$ ($n = 9$, $\ell = 5$, $|\SYT| = 120$).}
\begin{tabular}{lrr}
\toprule
Stratum & Count & Explanation \\
\midrule
$\ker(\Om)$ basis & $105$ & col-zero (Lemma~\ref{lem:col-zero}) \\
$\ker(\Omstar)$ basis & $48$ & row-zero (Lemma~\ref{lem:row-zero}) \\
Intersection & $45$ & both \\
Union & $108$ & all zero $\pT$ \\
Neither & $12$ & all nonzero $\pT$ \\
\midrule
Of the $105$ column-zero: pair $(1,2)$ SC & $70$ & descent-$1$ trivial \\
\quad\ldots\ deeper col cancellation & $35$ &  \\
Of the $48$ row-zero: pair $(5,6)$ SC & $26$ & Prop.~\ref{prop:row-suff} \\
\quad\ldots\ deeper row cancellation & $22$ &  \\
\midrule
$\ker(\Omstar) \setminus \ker(\Om)$ & $3$ & interior zeros $\{2, 5, 20\}$ \\
\quad among these, D-class & $2$ & $T_a, T_b$ (today's targets) \\
\quad among these, SC & $1$ & $T_0$ (May~17 smoking gun) \\
\bottomrule
\end{tabular}
\end{table}

\subsection{Distribution of D-class zeros}

Of the $108$ zero diagonals:
\begin{itemize}[topsep=0.2em,itemsep=0.1em]
\item $26$ are SC at $(\ell, \ell+1)$ (\emph{closed} by Diagnostic~3 SC,
   \cite{leftmostPerSYT}).
\item $4$ are SR at $(\ell, \ell+1)$ (\emph{closed} by column-zero via
   pair-$(1,2)$ SC, since all $4$ also satisfy that condition; combined
   with descent-1 mechanics).
\item $78$ are D at $(\ell, \ell+1)$ (the ``D-class''):
   \begin{itemize}[label=$\circ$,topsep=0.2em,itemsep=0.1em]
   \item $76$ have $\Om v_T = 0$ (column-zero, of which $52$ have pair
      $(1, 2)$ SC, and $24$ have pair $(1, 2)$ SR or D but are killed by
      deeper iterated rightmost-factor cancellation in $\Rp_7 \Rp_8 \Rp_9$).
   \item $\mathbf{2}$ have $\Om v_T \neq 0$ but row $T$ of $\Om$ is zero:
      these are the ``interior D-class zeros'' $T_a, T_b$, closed by
      Lemma~\ref{lem:row-zero}.
   \end{itemize}
\end{itemize}

The two interior D-class zeros were the explicit open problem at the head
of the May-17 SC paper's discussion section~\cite[\S5.3]{leftmostPerSYT};
Theorem~\ref{thm:main} together with Lemma~\ref{lem:row-zero} closes them.

\section{Open: the structural characterisation of $\ker(\Omstar)$}\label{sec:open}

The sufficient direction of Theorem~\ref{thm:main} is now proved
(Section~\ref{sec:sufficient}). The full equivalence is verified
computationally at $\lambda = (3,3,1,1,1)$ but is not yet proved
structurally. The remaining gap is a clean combinatorial description of
$\ker(\Omstar)$ on $\Vlam$.

\begin{conjecture}[Stratification of $\ker(\Omstar)$]\label{conj:kerstar}
For $\lambda = (\hatl, 1^q)$ in the standing hypothesis, the kernel
$\ker(\Omstar) \subseteq \Vlam$ stratifies according to ``leftmost-factor
elimination'' in $\Lp_n \cdots \Lp_{\ell+1}$: a basis vector $v_T$
satisfies $\Omstar v_T = 0$ iff there exists $k \in \{\ell+1, \ldots, n\}$
such that the partial product $\Lp_k \Lp_{k-1} \cdots \Lp_{\ell+1}$
applied to $v_T$ lies in
$\ker(\Lp_{k+1})$ in $\Vlam$ (with the convention $\Lp_{n+1} = I$).
\end{conjecture}

The $k = \ell+1$ case is Proposition~\ref{prop:row-suff}: $\Lp_{\ell+1}
v_T = 0$. The $k > \ell+1$ cases would explain the additional $22$
row-zero basis vectors at $(3,3,1,1,1)$, including the $T_a, T_b$ cases.
A direct trace of $\Lp_n \Lp_{n-1} \Lp_{n-2} \cdot \Lp_{\ell+1} v_{T_a}$
should pin down the precise cancellation (a $2$-block traversal through
pair $(5, 6)$ D, followed by a single $S_4$ killing the off-diagonal
contribution at the $S_4 v_{s_5 T_a}$ step). We leave this as the
strongest near-target generalisation.

Analogously, the column-zero side admits the conjecture:

\begin{conjecture}[Stratification of $\ker(\Om)$]\label{conj:keromega}
For $\lambda = (\hatl, 1^q)$ in the standing hypothesis, $\ker(\Om)$
stratifies by ``rightmost-factor elimination'' in $\Rp_{\ell+1} \cdots
\Rp_n$ symmetrically: $\Om v_T = 0$ iff there exists $k$ with
$\Rp_k v_T \in \ker(\Rp_{k+1} \cdots \Rp_{n-1} \cdot \Rp_n)$.
\end{conjecture}

The two conjectures together would give a fully combinatorial
characterisation of zero diagonals, refining the survey
diagnostics~\cite[\S7]{survey} to the algorithmic
``trace the cancellation chain on each side''.

\begin{remark}[$\tau \ge 1$ generalisation]
For $\hatl$ with $\tau(\hatl) \ge 1$ in the Pillar-1 regime, the rank
$\dim B = 0$ (by the trace-vanishing forward theorem), so $B = 0$, every
$\pT = 0$, and every basis vector is in $\ker(\Om) \cap \ker(\Omstar)$.
The two-sided dichotomy is vacuously true. The interesting regime is
$\tau = 0$ with $\dim B > 0$, where Theorem~\ref{thm:main} provides a
non-trivial decomposition.
\end{remark}

\begin{remark}[Why ``row-zero'' is the right mirror of Diagnostic~3]
Diagnostic~3 in its survey form~\cite[\S7]{survey} packages two facts:
(a) the leftmost-factor SC theorem~\cite[Thm.~1]{leftmostPerSYT}
(``$v_T \in \Eminus_\ell \Rightarrow \pT = 0$''), and
(b) the SR companion~\cite[Thm.~1]{srCompanion}
(``$v_T \in \Eplus_j$ for $j \le \tau$ $\Rightarrow \pT = 0$''). Both
exploit \emph{containment of one side} of $\Vlam = \Eplus_i \oplus
\Eminus_i$ in either $B$ (for (a)) or $\Vlam / B$ (for (b)).
Lemma~\ref{lem:row-zero} is the precise mirror: instead of \emph{one
specific} eigenspace constraint, it allows \emph{any} $v_T$
annihilated by the full anti-involution image. The $26$ SC tableaux are
the subset of $\ker(\Omstar)$ explained by the single-step
$S_\ell$-rightmost-factor; the additional $22$ in $\ker(\Omstar)$ at
$(3,3,1,1,1)$ are the deeper cancellations that
Conjecture~\ref{conj:kerstar} would explain.
\end{remark}

\begin{thebibliography}{9}
\bibitem{leftmostPerSYT}
   Clio, \emph{Per-SYT vanishing at the leftmost factor: $\pT = 0$ whenever
   pair $(\ell, \ell+1)$ is same-column},
   17 May 2026.
   \texttt{\textasciitilde/projects/proofs/2026-05-17-leftmost-factor-per-syt-vanishing.tex}.

\bibitem{srCompanion}
   Clio, \emph{A per-SYT diagonal vanishing companion: same-row pairs
   and Pillar~1}, 18 May 2026.
   \texttt{\textasciitilde/projects/proofs/2026-05-18-sr-companion.tex}.

\bibitem{pillar1}
   Clio, \emph{Extended sign-kill for general multi-row $\hatl$: the
   meta-theorem at arbitrary $\elh \ge 2$}, 15 May 2026.
   \texttt{\textasciitilde/projects/proofs/2026-05-15-multi-row-sign-kill-meta.tex}.

\bibitem{survey}
   Clio, \emph{The trace-vanishing dichotomy at
   $\lambda = (\hatl, 1^q)$: survey}, 17 May 2026.
   \texttt{\textasciitilde/projects/proofs/2026-05-17-survey-trace-vanishing-dichotomy.tex}.

\bibitem{verifyConv}
   Clio, \emph{Converse trace-vanishing: per-SYT positivity and
   column-descent-free SYT existence}, 17 May 2026.
   \texttt{\textasciitilde/projects/proofs/2026-05-17-converse-per-syt-positivity.tex}.

\bibitem{Murphy}
   G.~E.~Murphy, \emph{A new construction of Young's seminormal
   representation of the symmetric group}, J.~Algebra \textbf{69} (1981)
   287--297. [Symmetric form of the seminormal basis; self-adjointness of
   $T_i$ in the $b\!=\!1$ normalisation.]
\end{thebibliography}

\end{document}
